\section{Objections and Replies}

\subsection*{If the Mass does not take us to Calvary, is it only symbolic?}

No. The alternative between physical time travel and empty symbolism is false.
Catholic sacramental causality is real precisely as sacramental. The whole
Christ is made truly, really, and substantially present under either species;
his one sacrifice is represented and offered; its fruit is applied; the Church
truly participates. None of these claims depends on the assembly acquiring the
historical coordinates of Jerusalem.

\subsection*{How could the Supper anticipate a sacrifice not yet completed?}

It did not make the completed historical death a past event. Christ then
truly offered his Body and Blood under sacramental signs while freely moving
toward their bloody immolation. His words identified the gifts with the Body
to be given and Blood to be poured out; his divine authority instituted the
sacramental mode; his obedient self-offering gave the act its unity.
``Anticipation'' preserves this relation without pretending to solve it
through chronology.

\subsection*{Does ``once for all'' make later offering impossible?}

It makes another immolation impossible and unnecessary. Hebrews opposes
Christ's single effective offering to repeated sacrifices unable to perfect
the worshipper. The Mass does not put Christ to death again. It sacramentally
offers the same living Christ whose completed self-offering has permanent
efficacy. Later application and participation do not imply deficiency in the
Cross any more than later baptisms imply another death and resurrection of
Christ.

\subsection*{Does the heavenly liturgy make every Mass timeless?}

The worship of the risen Lamb and Christ's enduring intercession establish a
real heavenly dimension (Heb 7:25; 9:24; Rev 5). ``Timeless,'' however, can
confuse divine eternity, angelic or heavenly duration, Christ's glorified human
life, and earthly liturgical succession. Each Mass occurs at a real created
time and place. Its communion with heavenly worship does not abolish those
coordinates.

\subsection*{May a preacher still use the forward/backward image?}

Only with immediate qualification, and usually a better phrase is available.
An image may be rhetorically memorable while catechetically expensive. If it
causes hearers to imagine renewed suffering, literal relocation, a merely
mental reconstruction, or a Eucharist facing only backward, it has failed.
``The Supper sacramentally anticipates the Cross; the Mass sacramentally makes
the one Paschal sacrifice present'' is less cinematic and more faithful.

